% Work-MAT | 2. rocnik | hodina 14
\documentclass[11pt,a4paper]{article}
\input{../../sablony/pripravy/priprava.sty}
\usepackage{tikz}
\usetikzlibrary{arrows.meta}
\nastavPripravu{2. ročník}{Exponenciální rovnice}{14}{Exponenciální rovnice II: společný základ}
\begin{document}
\poznamka{Nácvik pravidel $a^{p}a^{q}=a^{p+q}$ a $(a^p)^q=a^{pq}$. Orientačně 6 min připomenutí, 12 min převody základů, 12 min obtížnější rovnice, 10 min samostatná práce, 5 min shrnutí.}
\blok{Nejdřív sjednotit základy}
Čísla $4,8,16,32$ jsou mocniny dvojky a $9,27,81$ mocniny trojky. Proto rovnice $4^x=8$ není vhodná pro přímé porovnání exponentů, dokud nepřepíšeme obě strany do stejného základu: $(2^2)^x=2^3$.
\begin{zadaniBox}[title={Věta - Pravidla pro mocniny}]
Pro $a>0$ a reálná $p,q$ platí $(a^p)^q=a^{pq}$, $a^pa^q=a^{p+q}$ a $a^p/a^q=a^{p-q}$. Mocniny s různými základy lze porovnat podle exponentů teprve po sjednocení na stejný základ $a\ne1$.
\end{zadaniBox}
\begin{prikladReseniBox}[title={Příklad 1}]
Vyřeš $9^{x-1}=27$.
\tcblower
\[\begin{aligned}(3^2)^{x-1}&=3^3\\3^{2x-2}&=3^3\\2x-2&=3\\x&=\frac52.\end{aligned}\]
$K=\{5/2\}$.
\end{prikladReseniBox}
\begin{prikladReseniBox}[title={Příklad 2}]
Vyřeš $8^{x+1}=4^{2x-1}$.
\tcblower
\[\begin{aligned}(2^3)^{x+1}&=(2^2)^{2x-1}\\2^{3x+3}&=2^{4x-2}\\3x+3&=4x-2\\x&=5.\end{aligned}\]
Zkouška v exponentu: vlevo $3\cdot6=18$, vpravo $2\cdot9=18$; obě strany jsou $2^{18}$.
\end{prikladReseniBox}
\begin{prikladReseniBox}[title={Příklad 3}]
Vyřeš $\left(\frac14\right)^{x-1}=8$.
\tcblower
\[\begin{aligned}(2^{-2})^{x-1}&=2^3\\2^{-2x+2}&=2^3\\-2x+2&=3\\x&=-\frac12.\end{aligned}\]
$K=\{-1/2\}$.
\end{prikladReseniBox}
\poznamka{Nezaměňuj $2^{x+2}$ za $2^x+2^2$. U převodu základu vždy nejdříve přepiš základ i exponent do závorky, například $(2^3)^{x-1}$.}
\begin{zadaniReseniBox}[title={Samostatná práce 1}]
\begin{enumerate}[label=\alph*)]
\item $25^{x-1}=125$;
\item $16^x=8^{x+2}$;
\item $\left(\frac19\right)^{x}=27$.
\end{enumerate}
\tcblower
\begin{enumerate}[label=\alph*)]
\item $5^{2x-2}=5^3$, tedy $x=\frac52$.
\item $2^{4x}=2^{3x+6}$, tedy $x=6$.
\item $3^{-2x}=3^3$, tedy $x=-\frac32$.
\end{enumerate}
\end{zadaniReseniBox}
\begin{zadaniReseniBox}[title={Rozšíření 1}]
Vyřeš $2^{x+2}\cdot4^{x-1}=32$.
\tcblower
$2^{x+2}\cdot2^{2x-2}=2^{3x}=2^5$, tedy $3x=5$ a $x=\frac53$.
\end{zadaniReseniBox}
\begin{zadaniReseniBox}[title={Domácí úkol}]
\begin{enumerate}[label=\alph*)]
\item $27^x=9^{x+1}$.
\item $\left(\frac18\right)^{x}=4^{x-1}$.
\item Vysvětli, proč pro $a>0$ platí $a^0=1$.
\end{enumerate}
\tcblower
\begin{enumerate}[label=\alph*)]
\item $3^{3x}=3^{2x+2}$, tedy $x=2$.
\item $2^{-3x}=2^{2x-2}$, tedy $x=\frac25$.
\item Z podílu $a^m/a^m=a^{m-m}=a^0=1$ pro $a>0$.
\end{enumerate}
\end{zadaniReseniBox}
\end{document}
